\documentclass[11pt,a4paper]{article}
\usepackage[utf8]{inputenc}
\usepackage[T1]{fontenc}
\usepackage{lmodern}
\usepackage[portuguese]{babel}
\usepackage[a4paper,margin=2.5cm]{geometry}
\usepackage{xcolor}
\definecolor{TextEspresso}{HTML}{4A4238}
\definecolor{NudeTaupe}{HTML}{BFAFA6}
\definecolor{SoftBlush}{HTML}{D9C5B2}
\usepackage{titlesec}
\usepackage{enumitem}
\usepackage{listings}
\usepackage{hyperref}
\usepackage{amsmath}
\hypersetup{colorlinks=true, linkcolor=TextEspresso, urlcolor=TextEspresso}

\titleformat{\section}{\color{TextEspresso}\Large\bfseries}{\thesection}{0.5em}{}[{\color{NudeTaupe}\titlerule[0.8pt]}]
\titleformat{\subsection}{\color{TextEspresso}\large\bfseries}{\thesubsection}{0.5em}{}

\lstset{
  basicstyle=\ttfamily\small,
  breaklines=true,
  frame=single,
  rulecolor=\color{NudeTaupe},
  columns=fullflexible,
  showstringspaces=false
}

\newcommand{\guiaotitle}[2]{%
  \begin{center}
  {\Large\bfseries\color{TextEspresso} #1}\\[0.3cm]
  {\normalsize\color{NudeTaupe} #2}\\[0.2cm]
  {\color{NudeTaupe}\rule{4cm}{0.5pt}}
  \end{center}
  \vspace{0.5cm}
}

\begin{document}
\guiaotitle{Guião 02 -- Análise de Arquitetura e Disponibilidade de Datacenters}{Infraestruturas de Computação na Nuvem, Aula 02}

\section{Objetivos}
\begin{itemize}[itemsep=0.2cm]
  \item Calcular a disponibilidade combinada de sistemas com componentes em série e em paralelo (redundância).
  \item Classificar um cenário de datacenter segundo os Tiers do Uptime Institute.
  \item Propor melhorias de redundância a uma arquitetura existente.
  \item Observar, de forma prática, indícios de uma topologia de rede real com \texttt{ping} e \texttt{traceroute}.
  \item Analisar, no próprio servidor Proxmox da UC, um caso real de partilha de recursos e de single point of failure.
\end{itemize}

\section{Pré-requisitos e Material}
\begin{itemize}[itemsep=0.2cm]
  \item Conhecimento dos conceitos de Tier, redundância N/N+1/2N e "noves" de disponibilidade, apresentados na Aula 02.
  \item Calculadora ou folha de cálculo para os exercícios numéricos.
  \item Acesso a um terminal com os comandos \texttt{ping} e \texttt{traceroute} (ou \texttt{tracert} no Windows) disponíveis; a VM criada no Guião 01 pode ser usada para este efeito.
  \item Acesso à interface web do Proxmox VE do servidor central, tal como usado no Guião 01.
\end{itemize}

\section{Enquadramento Teórico Breve}
A disponibilidade de um sistema, $A$, representa a fração de tempo em que este se encontra operacional, e relaciona-se com o \emph{Mean Time Between Failures} (MTBF) e o \emph{Mean Time To Repair} (MTTR) através de:
\[
A = \frac{\text{MTBF}}{\text{MTBF} + \text{MTTR}}
\]

Para um conjunto de componentes independentes ligados em \textbf{série} (todos têm de funcionar para o sistema funcionar), a disponibilidade combinada é o produto das disponibilidades individuais:
\[
A_{\text{série}} = A_1 \times A_2 \times \cdots \times A_n
\]

Para um conjunto de componentes redundantes ligados em \textbf{paralelo} (basta um funcionar), a disponibilidade combinada é:
\[
A_{\text{paralelo}} = 1 - \big[(1-A_1)(1-A_2)\cdots(1-A_n)\big]
\]

O downtime anual associado a uma disponibilidade $A$ obtém-se por $(1-A) \times 365{,}25 \times 24 \times 60$ minutos.

\section{Exercícios}
\begin{enumerate}[label=\arabic*.]
  \item \textbf{Sistema em série.} Um serviço depende de três componentes ligados em série: o link de rede (disponibilidade 99,95\%), o servidor de aplicação (99,9\%) e a base de dados (99,99\%). Calcule a disponibilidade combinada do sistema e o downtime anual esperado, em horas.
  \vspace{2cm}

  \item \textbf{Sistema em paralelo.} A alimentação elétrica do datacenter é assegurada por duas fontes independentes (rede elétrica e gerador de emergência), cada uma com disponibilidade individual de 99,5\%. Calcule a disponibilidade combinada do sistema de alimentação e compare com a disponibilidade de uma única fonte.
  \vspace{2cm}

  \item \textbf{Sistema misto.} Um serviço depende de um link de rede (99,9\%) em série com um par redundante de servidores de aplicação (99\% cada, em paralelo). Calcule a disponibilidade final do sistema completo.
  \vspace{2cm}

  \item \textbf{Classificação Tier.} Um datacenter possui uma única via de alimentação elétrica e um único sistema de arrefecimento, mas com um gerador de reserva (N+1) para a alimentação. A manutenção de qualquer componente obriga à interrupção do serviço. Classifique este cenário segundo os Tiers do Uptime Institute e justifique a resposta.
  \vspace{2cm}

  \item \textbf{Proposta de melhoria.} Com base no cenário do exercício anterior, proponha duas alterações concretas à arquitetura que permitiriam evoluir para um Tier III. Justifique o impacto de cada alteração na disponibilidade global.
  \vspace{2cm}

  \item \textbf{Observação de topologia de rede.} A partir de um terminal, execute os comandos seguintes para um destino à sua escolha (por exemplo, \texttt{www.ua.pt}) e registe os resultados:
\begin{lstlisting}
$ ping -c 5 www.ua.pt
$ traceroute www.ua.pt
\end{lstlisting}
  Analise o resultado do \texttt{traceroute}: quantos saltos (hops) são percorridos até ao destino? Existe alguma variação significativa de latência entre saltos consecutivos? O que pode isso indicar sobre a topologia de rede atravessada?
  \vspace{1cm}

  \item \textbf{Análise do servidor central da UC.} Aceda à interface web do Proxmox VE (a mesma do Guião 01) e, no separador \emph{Summary} do node, registe os valores observados:
\begin{lstlisting}
- Uptime do servidor
- Numero de CPUs (cores/threads) e utilizacao atual
- Memoria RAM total e em uso
- Espaco de armazenamento total e disponivel
\end{lstlisting}
  Com base nestes dados e no que observa da infraestrutura disponível, responda:
  \begin{enumerate}[label=\alph*)]
    \item Que Tier do Uptime Institute atribuiria a esta infraestrutura? Justifique identificando concretamente que elementos de redundância existem ou faltam.
    \item Assumindo que o servidor tem uma única fonte de alimentação e uma única ligação de rede, e considerando uma disponibilidade individual de 99\% para cada um destes dois componentes em série, qual seria a disponibilidade combinada estimada e o downtime anual correspondente?
    \item Identifique pelo menos dois \emph{single points of failure} nesta infraestrutura e, para cada um, indique que alteração concreta os eliminaria.
  \end{enumerate}
  \vspace{1cm}

  \item \textbf{Latência dentro vs fora do datacenter.} A partir da sua VM (criada no Guião 01), compare a latência para um destino \emph{dentro} da infraestrutura local com a latência para um destino externo:
\begin{lstlisting}
$ ping -c 5 <IP-do-servidor-Proxmox>
$ ping -c 5 www.ua.pt
\end{lstlisting}
  Compare os valores médios obtidos. Relacione a diferença observada com a distinção entre tráfego \emph{este-oeste} (dentro do datacenter) e \emph{norte-sul} (para o exterior) apresentada na aula teórica.
\end{enumerate}

\section{Questões de Verificação}
\begin{itemize}[itemsep=0.2cm]
  \item Porque motivo a disponibilidade combinada de componentes em série é sempre inferior à do componente mais fraco da cadeia?
  \item Explique, em termos práticos, a diferença entre redundância N+1 e redundância 2N.
  \item Que vantagem tem a topologia leaf-spine sobre a three-tier relativamente à previsibilidade da latência observada num \texttt{traceroute}?
  \item Todas as VMs da turma correm sobre o mesmo servidor físico. Que característica essencial da computação na nuvem (Aula 01) isto ilustra, e que risco operacional introduz?
\end{itemize}

\end{document}
